\documentclass[12pt,a4paper]{article}

\usepackage[utf8]{inputenc}
\usepackage[T1]{fontenc}
\usepackage[french]{babel}

\usepackage{styles/couleurs}
\usepackage{styles/mise_en_page}
\usepackage{styles/boites}
\usepackage{styles/maths}
\usepackage{styles/figures}

\usepackage{amsmath}
\usepackage{amssymb}
\usepackage{enumitem}
\usepackage{array}
\usepackage{tabularx}
\usepackage{graphicx}
\usepackage{tikz}
\usepackage{xparse}
\usetikzlibrary{calc,angles,quotes}

\setlist[itemize]{leftmargin=1.4em,itemsep=0.15em,topsep=0.2em}
\setlist[enumerate]{leftmargin=*,itemsep=0.22em,topsep=0.22em}

\newcommand{\titresequence}[1]{%
\begin{center}
  {\sffamily\LARGE\bfseries\color{TitreC}Séquence 2 -- Théorème de Pythagore}\\[0.2em]
  {\sffamily\large\bfseries\color{SavoirC}#1}
\end{center}}

\newcommand{\lignepointillee}{\noindent\dotfill\par\vspace{0.45em}}
\newcommand{\lignesreponse}[1]{%
\begin{tcolorbox}[
  enhanced,colback=white,colframe=NeutreC!40,boxrule=0.5pt,arc=1.5mm,
  left=3mm,right=3mm,top=2mm,bottom=1mm,before skip=2pt,after skip=2pt]
  \foreach \i in {1,...,#1}{\lignepointillee}
\end{tcolorbox}}

\begin{document}
\titresequence{Activité 3 -- Introduction rapide à la racine carrée}

On veut construire un carré dont l'aire est $50$ cm$^2$.
On cherche la longueur de son côté.

\begin{enumerate}
  \item Complète :
  \[
  7^2=\dotfill
  \qquad\text{et}\qquad
  8^2=\dotfill
  \]
  Entre quels deux entiers se trouve la longueur cherchée ?
  \lignesreponse{3}

  \item Calcule :
  \[
  7{,}0^2=\dotfill
  \qquad
  7{,}1^2=\dotfill
  \qquad
  7{,}2^2=\dotfill
  \]
  Donne un encadrement au dixième de la longueur cherchée.
  \lignesreponse{3}

  \item La longueur cherchée se note $\sqrt{50}$.
  Complète :
  \begin{quote}
  La racine carrée de $50$ est le nombre positif qui, multiplié par lui-même, donne \dotfill
  \end{quote}

  \item À la calculatrice, donne une valeur approchée au dixième de $\sqrt{50}$.
  \lignesreponse{2}
\end{enumerate}

\definition{La \textbf{racine carrée} d'un nombre positif $a$ est le nombre positif dont le carré vaut $a$.
On la note $\sqrt{a}$.}

\exemple{
\[
\sqrt{25}=5
\qquad
\sqrt{49}=7
\qquad
\sqrt{81}=9.
\]}

\section*{Retour à Pythagore}

Le triangle $RST$ est rectangle en $R$, avec $RS=2$ cm et $RT=7$ cm.

\begin{enumerate}
  \item Calcule $ST^2$.
  \item Donne la valeur exacte de $ST$.
  \item Donne une valeur approchée au dixième de $ST$.
\end{enumerate}
\lignesreponse{7}
\end{document}
